Several other papers have studied the estimation of frequentist variability of
posterior expectations.  \citet{efron:2015:frequentist} provides results similar
to ours, but is limited to the special case of exponential families. The
frequentist sampling variability of Bayesian posterior expectations under
misspecification is studied in \citet{kleijn:2012:bvm} from a theoretical point
of view, motivating the form of the sandwich covariance.
\Citet{huggins:2024:bayesbagparameter} derive a BCLT for the ``BayesBag''
posterior which combines both posterior and bootstrap uncertainty; we conjecture
that a minor modification of their proof would prove the consistency of the
bootstrap alone as an estimator of the limiting frequentist variance of a
posterior mean. See also \citet{huggins:2022:bayesbag} for a closely related
BCLT result in the context of model selection, as well as the same authors'
earlier work on arxiv, \citet{huggins:2020:bayesbag}.  Although the IJ
covariance does not provide bootstrap samples for use in BayesBag posteriors
directly, the IJ covariance could be used to estimate the bootstrap covariance
term that appears in the law of total covariance decomposition of the BayesBag
covariance (see the final, unnumbered equation in Section 2.3 of
\citet{huggins:2024:bayesbagparameter}). Both \citet{vehtari:2024:psis} and
\citet{lee:2017:frequentistbayesse} consider importance sampling estimates to
the effect of dropping datapoints, with \citet{lee:2017:frequentistbayesse}
explicitly considering estimation of frequentist variability of the posterior
under an assumption of correct specification. However, importance sampling
estimates of frequentist variability can be expected to have very low sampling
efficiency even in moderate dimensional problems, since frequentist re--sampling
is expected to produce changes of the posterior mean on the same order as the
posterior standard deviation (see Example 9.2 of \citet{owen:2014:mcbook}).
Furthermore, misspecification can result in frequentist sampling variances that
are different than the posterior variance, which can further degrade importance
sampling efficiency as a function of dimension (see Example 9.3 of
\citet{owen:2014:mcbook}).

Our key result, \thmref{bayes_clt_main} is distinguished from previous BCLT
results primarily by tracking the data dependence in the residual of the series
expansion, though we also differ from some classical work by admitting
misspecification.  Our method of proof hews closely to the classical BCLT proofs
found in  Chapter 10 of \citet{van:2000:asymptotic}, Theorem 1.4.2 of
\citet{ghosh:2003:bayesian}, and \citet{lehman:1983:pointestimation}, all of
which study the asymptotic regime for a single posterior quantity under the
assumption of correct specification.  A series expansion for posterior
expectations similar to \thmref{bayes_clt_main} is established by
\citet{kass:1990:posteriorexpansions}, under stricter differentiability
conditions than ours, and only for a single posterior quantity of interest,
which is (implicitly) assumed not to depend on the dataset. A number of recent
papers have, like us, studied the asymptotic behavior of finite--dimensional
Bayesian posteriors under misspecification, though without explicitly
controlling data dependence
\citep{kleijn:2006:misspecification,miller:2021:bclt,kasprzak:2022:laplace}.

Our work is an application of the classical IJ estimator, which has been
extensively studied in the frequentist literature, but which to the best of the
authors' knowledge has not been applied to Bayesian posterior expectations (see,
e.g., \citet{jaeckel:1972:infinitesimal}, \citet{efron:1982:jackknife},
\citet{shao:2012:jackknife}, \citet{efron:2014:estimation},
\citet{wager:2014:confidence},
\citet{peng:2021:biasconsistencyalternativeperspectives},
\citet{ghosal:2022:infinitesimaljackknifecombinationsmodels}).
As we discuss more explicitly in \secref{data_weights}, the IJ covariance can be
viewed as the first term in a von Mises expansion of the Bayesian posterior
(\citet{mises:1947:asymptotic}, \citet{reeds:1976:thesis},
\citet{serfling:1980:approximation}), though we prove consistency directly
rather than via a functional series expansion.
\Citet{giordano:2024:oldbayesij} prove consistency of the IJ for posterior
expectations using von Mises expansions directly under stronger theoretical
assumptions, using results from the present paper.

As we discuss in the introduction, IJ estimation can be thought of as a
particular measure of local data robustness --- specifically, robustness to
random resampling of the data. We make this connection even more explicit in
\secref{data_weights} below. Viewing IJ estimators as local data robustness
connects our work to a large literature on ``local case influence,'' which
refers to a broad class of techniques that using derivatives to approximate the
effect of a datapoint or datapoints on a statistic of interest.  Local case
influence is an idea with a long history in frequentist and Bayesian statistics.
The literature is too vast to review systematically here, but some notable
references include
\citet{belsley:1980:regression,cook:1977:detectionofinfluential,
cook:1982:residualsandinfluenceinregression, cook:1986:assessment,
kempthorne:1986:decision,kass:1989:approximatebayesiansensitivity,
carlin:1991:expectedutilityinfluence,
zhu:2007:perturbation,koh:2017:blackboxinfluence,
thomas:2018:reconcilingcurvatureandis}.
The relationship between posterior covariances and derivatives has been noted
numerous times in the Bayesian robustness literature (e.g.
\citet{diaconis:1986:bayesconsistency}, \citet{gustafson:1996:localposterior},
\citet{basu:1996:local}, \citet{efron:2015:frequentist},
\citet{giordano:2018:covariances}). As we note in \secref{bayes_clt},
consistency of the IJ requires a kind of uniform guarantee over a large number
of local approximations, but most of the local case influence work focuses on a
datapoint or a small set of datapoints. Notable exceptions include
\citet{giordano:2019:swiss} in the frequentist literature and
\citet{zhu:2012:bayesiancaseinfluence} in the Bayesian literature, though we
note that Theorem 2 of \citet{zhu:2012:bayesiancaseinfluence} is stated in terms
of a fixed set of dropped datapoints.

Following the posting of the present work on the arXiv, a number of interesting
subsequent papers have proposed extensions and variants and applications of the
Bayesian IJ. \Citet{nguyen:2024:sensitivitymcmcbasedanalysessmalldata} applies
the IJ to the adversarial data--dropping work of
\citet{giordano:2026:amip1,giordano:2026:amip2} in MCMC settings.
\Citet{iba:2023:wkernel} connects the IJ covariance to kernel methods, and
\citet{iba:2023:posteriorweighted} apply ideas closely related to the Bayesian
IJ covariance to information criteria and model selection.  Two papers,
\textcite{ji:2025:ValidStandardErrors2025,luo:2026:RobustStandardErrors2026},
have systematically evaluated the Bayesian IJ covariance for use in the social
sciences in the presence of model misspecification. We hope that our approach
continues to motivate and inform connections between the frequentist and
Bayesian case influence literature.




